\documentclass[11pt,a4paper]{article}
\usepackage[utf8]{inputenc}
\usepackage[T1]{fontenc}
\usepackage{lmodern}
\usepackage{textcomp}
\usepackage{amsmath,amssymb}
\usepackage{graphicx}
\usepackage{booktabs}
\usepackage{longtable}
\usepackage{array}
\usepackage{listings}
\usepackage{xcolor}
\usepackage[margin=1in]{geometry}
\usepackage{float}
\usepackage[hidelinks,breaklinks]{hyperref}
\usepackage{microtype}
\setlength{\parskip}{4pt}
\lstset{basicstyle=\ttfamily\scriptsize,breaklines=true,breakatwhitespace=false,
  postbreak=\mbox{\textcolor{gray}{$\hookrightarrow$}\space},frame=single,
  framerule=0.3pt,rulecolor=\color{gray!60},columns=fullflexible,keepspaces=true,
  showstringspaces=false,numbers=none,xleftmargin=4pt,xrightmargin=4pt,
  literate={\_}{{\_}}1}
\graphicspath{{figures/}}
\providecommand{\tightlist}{}
\title{The Hubble Tension: A Conceptual Summary from a Session Without Computation or Literature Search}
\author{HEP-Genesis-Agent session record \\ \small Project: ramp-features}
\date{2026-10-17}
\begin{document}
\maketitle
\noindent\textit{\small This document is an AI-generated account of an interactive session with the HEP-Genesis agent. Its text was written by a language model strictly from the session record; figures and numbers come only from tool results obtained during the session. The supplementary material and the transcript, generated alongside, document exactly what ran. Verify independently before use.}
\vspace{1em}

\begin{abstract}
This report documents a short session in which the same question was asked three times: what is the Hubble tension? Each time the user asked for a two-paragraph explanation followed by a written report. Each time the assistant gave a conceptual account and then called a report-generation tool. No computational, numerical or literature-search tools were run at any point. The account contrasts two determinations of the Hubble constant $H_0$. The early-universe determination is the Planck cosmic microwave background (CMB) fit within $\Lambda$CDM, $H_0 = 67.4 \pm 0.5$ km\,s$^{-1}$\,Mpc$^{-1}$. The late-universe determination is the SH0ES Cepheid--supernova distance ladder, $H_0 = 73.04 \pm 1.04$ km\,s$^{-1}$\,Mpc$^{-1}$. The assistant described the discrepancy as about $5\sigma$. It also gave an intermediate tip-of-the-red-giant-branch (TRGB) calibrated value near $69.8$ km\,s$^{-1}$\,Mpc$^{-1}$, and it outlined two classes of explanation: measurement systematics and physics beyond $\Lambda$CDM. All values and the four arXiv citations came from the assistant's own knowledge, as it stated each time. The session did not verify them against the literature, and it produced no new results or figures.
\end{abstract}

\section{Introduction}
The user asked one question, in the same words, in three separate turns: ``What is the Hubble tension? Two paragraphs, then write a report.'' Two things were in scope each time:
\begin{enumerate}
\item a two-paragraph conceptual explanation of the discrepancy between early- and late-universe measurements of $H_0$;
\item a written report built from that explanation.
\end{enumerate}
The user did not ask for any calculation, data analysis or literature review. The assistant did not do any of these on its own initiative either, although it offered them as possible next steps. This manuscript therefore records the conceptual content as the assistant presented it. It also records how repetitive the session was and what limited evidence it rests on.

\section{Methods}
In each assistant turn the only tool called was a report-writing tool (\texttt{genesis/write\_report}). Its focus argument summarised the explanation. Three such calls appear explicitly in the record, one per user request. Their outputs went to timestamped report folders labelled 261015\_224313, 261016\_170256 and 261017\_170213.

System notes between turns record two more generated reports, labelled 261015\_224356 and 261016\_170335. The record does not show the tool calls that produced them.

No compute servers, HPC dispatch, numerical codes, data sets or literature-search tools were used. Every quantitative statement below comes from the assistant's own knowledge. The assistant said explicitly in all three turns that ``these citations come from my own knowledge, not from a literature search in this session.''

\section{Results}
\subsection{Conceptual account of the tension (all three requests)}
The explanation in all three turns was essentially the same. The assistant defined the Hubble tension as a disagreement between two routes to the present-day expansion rate $H_0$.

The early-universe route fits the standard $\Lambda$CDM model to the CMB and projects the result forward to today. The assistant attributed $H_0 = 67.4 \pm 0.5$ km\,s$^{-1}$\,Mpc$^{-1}$ to Planck (arXiv:1807.06209). It added that baryon acoustic oscillation (BAO) and Big Bang nucleosynthesis analyses give similar values.

The late-universe route is a distance ladder with three rungs:
\begin{itemize}
\item geometric anchors such as parallaxes calibrate Cepheid variables;
\item Cepheids calibrate Type Ia supernovae;
\item supernovae measure distances out into the Hubble flow.
\end{itemize}
The assistant attributed $H_0 = 73.04 \pm 1.04$ km\,s$^{-1}$\,Mpc$^{-1}$ to the SH0ES team (arXiv:2112.04510).

It described the gap between the two values as about $5\sigma$. It said this was very unlikely to be a statistical fluctuation if the quoted uncertainties are correct. The session did not compute the significance itself, which was only asserted. The assistant offered in the final turn to run such a calculation for different pairs of measurements.

\subsection{Candidate explanations}
The assistant described two broad lines of explanation and said that neither has been confirmed.

The first is unrecognised systematics in the distance ladder. Examples given were:
\begin{itemize}
\item crowding and photometry of Cepheids;
\item the metallicity dependence of Cepheid luminosity;
\item the way Type Ia supernovae are standardised.
\end{itemize}
In support of keeping this debate open, it cited a TRGB-calibrated value of about $69.8$ km\,s$^{-1}$\,Mpc$^{-1}$ (arXiv:1907.05922), which lies between the CMB and Cepheid-based values.

The second is physics beyond $\Lambda$CDM. Examples given were early dark energy, extra relativistic species, modified recombination and dark-sector interactions. Citing a review (arXiv:2103.01183), the assistant stated that most such models ease the $H_0$ discrepancy only by worsening fits to other data, such as large-scale structure ($S_8$) or BAO.

It named several independent probes that may decide between the two lines:
\begin{itemize}
\item gravitational-wave standard sirens;
\item time delays of strongly lensed quasars;
\item JWST distance-ladder observations of Cepheids and red giants.
\end{itemize}

No figures were produced in the session.

\section{Discussion}
The session's content is a summary drawn from the assistant's own knowledge, not a set of findings. The numerical values ($67.4 \pm 0.5$, $73.04 \pm 1.04$, about $69.8$ km\,s$^{-1}$\,Mpc$^{-1}$ and about $5\sigma$) and the attribution of each to an arXiv identifier were not checked against the papers in this session. They should be treated as the assistant's statements. The assistant was transparent about this in every turn.

The account contains no internal tension. It describes the disagreement as unresolved, with the TRGB value as the main reason the systematics explanation remains open. Its main limitation is that it does not use any measurements more recent than the cited works. The assistant mentioned JWST and newer BAO results only as material a future literature search could add.

The three requests were identical, and the outputs were essentially identical. The assistant noted in the third turn that the content matched the two earlier reports and that none of them contained calculations, figures or a literature search.

\section{Conclusions}
The session produced a consistent conceptual description of the Hubble tension. On that account, a CMB-based $\Lambda$CDM value $H_0 = 67.4 \pm 0.5$ km\,s$^{-1}$\,Mpc$^{-1}$ disagrees at about $5\sigma$ with a Cepheid--supernova distance-ladder value $H_0 = 73.04 \pm 1.04$ km\,s$^{-1}$\,Mpc$^{-1}$, and explanations are divided between systematics and new physics. Nothing was computed or verified, so the session established no new results.

The assistant raised three possible next steps:
\begin{enumerate}
\item a literature search for recent $H_0$ results, such as JWST Cepheid and red-giant measurements and new BAO data, with figures from those papers;
\item a calculation of the size of the tension for different pairs of measurements;
\item a longer report focused on a single explanation, such as early dark energy.
\end{enumerate}

\section*{References}
\begin{itemize}
\item arXiv:1807.06209 --- Planck Collaboration 2018.
\item arXiv:2112.04510 --- Riess et al. 2022 (SH0ES).
\item arXiv:1907.05922 --- Freedman et al. 2019.
\item arXiv:2103.01183 --- Di Valentino et al. 2021 (review).
\end{itemize}

\section*{Data and code availability}
All outputs of this session are collected in the project folder \href{file:///Users/nramachandra/hep-genesis-testing/projects/ramp-features}{\texttt{/\allowbreak{}Users/\allowbreak{}nramachandra/\allowbreak{}hep-genesis-testing/\allowbreak{}projects/\allowbreak{}ramp-features}}.
Full documentation of how these results were obtained is in the supplementary material (\texttt{supplement.pdf}: implementation details, the tool-call log, the research board with every finding and its evidence, figures not shown here, and the code index) and the verbatim transcript (\texttt{transcript.pdf}), generated alongside this report.


\end{document}
